%% ==============================================================================
%% content/04-dynamique-page.tex
%% Chapitre 4 : Maîtrise Dynamique de l'Empagement : KOMA-Script, Géométries,
%% En-têtes Avancés et Multi-Colonnes
%% ==============================================================================
\chapter{Maîtrise Dynamique de l'Empagement : KOMA-Script, Géométries et Multi-Colonnes}
\label{chap:dynamique-empagement}

\section{Fondements Typographiques et Empagement KOMA-Script}
\label{sec:fondements-typographiques}

La mise en page d'un document long (thèse doctorale, traité scientifique ou manuel d'ingénierie)
ne résulte jamais d'un choix esthétique arbitraire.
Elle obéit à des lois de physiologie visuelle et d'ergonomie cognitive théorisées par \textcite{tschichold1991form}
et rappelées avec force par \textcite{bringhurst2004elements}.

\subsection{Le canon de division typographique et la valeur \code{DIV}}
Le confort de lecture dépend étroitement de la longueur des lignes de texte :
l'\oe{}il humain fatigue rapidement si une ligne excède 65 à 75 caractères (soit 10 à 12 mots par ligne).
Pour respecter cette contrainte sans gaspiller le papier,
la classe KOMA-Script \code{scrbook} intègre l'algorithme de calcul \term{typearea} :

\begin{enumerate}
  \item La surface de la feuille est subdivisée en une grille de $N \times N$ rectangles identiques,
  déterminée par le coefficient \code{DIV} (dans notre classe, \code{DIV=12}).
  \item La zone de composition textuelle (\textit{Satzspiegel}) occupe $10 \times 10$ rectangles,
  tandis que les marges absorbent les subdivisions périphériques selon des proportions savantes.
  \item En reliure classique recto-verso (\code{print}), la marge intérieure (petit fond)
  mesure exactement la moitié de la marge extérieure (grand fond),
  car les deux pages ouvertes en regard fusionnent visuellement leurs petits fonds en une bande d'égale largeur.
\end{enumerate}

\subsection{Le défi des livrables modernes : au-delà de la géométrie unique}
Bien que cette rigueur typographique soit idéale pour la prose continue,
les travaux d'ingénierie contemporains se heurtent à une rigidité historique :
comment insérer une matrice comparative de 12 colonnes,
un tableau synoptique de synthèse de 180\,mm de large,
ou aménager des marges généreuses pour les révisions collaboratives,
sans rompre la cohérence globale du volume ?

La réponse conventionnelle consiste souvent à bricoler des valeurs globales dans \code{geometry},
au détriment de l'harmonie des 250 autres pages.
Dans notre architecture, nous combinons la puissance sémantique de KOMA-Script
avec le pilote géométrique de \code{geometry} grâce à l'option \code{\textbackslash KOMAoptions\{usegeometry=true\}}.
Cette symbiose permet de métamorphoser la géométrie d'une page à l'autre
avec une fluidité absolue.

\begin{boiteSynthese}{Les Trois Géométries Dynamiques de l'Architecture}
Notre infrastructure propose trois géométries spécialisées activables à la demande :
\begin{itemize}
  \item \textbf{\code{\textbackslash geometrieStandard} :}
  la géométrie canonique KOMA (\code{DIV=12}) préservant le ratio de lecture optimal.
  \item \textbf{\code{\textbackslash geometrieLarge} (ou environnement \code{WideGeometry}) :}
  déploie une zone utile élargie (marges de 15\,mm) pour les données denses.
  \item \textbf{\code{\textbackslash geometrieWord} / \code{\textbackslash geometrieAnnotation} :}
  aménage une marge latérale de révision de 25\,mm dédiée aux notes \code{\textbackslash todoH}.
\end{itemize}
\end{boiteSynthese}

% ==============================================================================
% 4.2 Démonstration en Direct de la Géométrie Large (WideGeometry)
% ==============================================================================
\section{Démonstration en Direct : Géométrie Pleine Page (\code{WideGeometry})}
\label{sec:demo-geometrie-large}

Pour illustrer cette plasticité, nous ouvrons ici même un bloc de géométrie large
à l'aide de l'environnement \code{WideGeometry}.
Remarquez le comportement sur les pages suivantes :
les marges latérales se réduisent instantanément à 15\,mm,
libérant une largeur d'empagement de plus de 180\,mm.
Simultanément, nous activons le style d'en-tête épuré \code{\textbackslash activerStyleFocus}
pour supprimer tout filet parasite et maximiser la concentration du lecteur.

\begin{WideGeometry}
\activerStyleFocus

\subsection{Espace élargi et style d'en-tête minimaliste \code{StyleFocus}}
Nous sommes actuellement dans l'environnement \code{WideGeometry}.
Observez l'en-tête supérieur : la ligne de séparation habituelle a disparu,
le titre courant de chapitre est masqué,
et seul un folio centré (\code{---\,\thepage\,---}) subsiste en bas de page.
Cette disposition est particulièrement recommandée lors de la présentation
de grands tableaux de bord de données ou de comparaisons exhaustives d'architectures.

Le \cref{tab:comparatif-architectures-large} exploite l'intégralité des 180\,mm disponibles
sans provoquer le moindre débordement (\textit{overfull hbox})
et sans devoir réduire la taille de police en corps microscopique.

\begin{table}[htbp]
  \centering
  \small
  \caption[Matrice comparative d'architectures documentaires en géométrie élargie]{Matrice comparative exhaustive des architectures documentaires (largeur utile : 180\,mm, environnement \code{WideGeometry})}
  \label{tab:comparatif-architectures-large}
  \begin{tabularx}{\linewidth}{>{\raggedright\arraybackslash}p{2.8cm} >{\raggedright\arraybackslash}X >{\raggedright\arraybackslash}X >{\raggedright\arraybackslash}X c}
    \toprule
    \textbf{Approche} & \textbf{Avantages majeurs} & \textbf{Inconvénients \& Risques} & \textbf{Cas d'usage privilégiés} & \textbf{Maturité} \\
    \midrule
    \textbf{Monolithique} &
    Fichier unique simple à appréhender pour les débutants. &
    Conflits Git constants, temps de compilation explosif, aucune réutilisabilité. &
    Articles courts (moins de 5 pages), devoirs ponctuels. &
    Obsolète \\
    \addlinespace
    \textbf{Modulaire Simple} &
    Découpage en chapitres via \code{\textbackslash input} ou \code{\textbackslash include}. &
    Frontière floue entre styles et données ; dépendances globales non contrôlées. &
    Mémoires de master, rapports techniques intermédiaires. &
    Courante \\
    \addlinespace
    \textbf{Découplée \code{mythesis}} &
    Isolation stricte (\code{.clo}, \code{config/}, \code{biblio/}, \code{listings/}). Profils écran/print instantanés. &
    Nécessite une rigueur initiale et l'utilisation de \code{latexmkrc}. &
    Thèses de doctorat, monographies, spécifications industrielles. &
    Optimale \\
    \addlinespace
    \textbf{Dynamique Hybride} &
    Alternance de géométries (\code{WideGeometry}, \code{AnnotationGeometry}) et calques interactifs. &
    Exige la maîtrise de KOMA-Script et de \code{scrlayer-scrpage}. &
    Ouvrages d'art, manuels de référence, revues à fort impact. &
    Avancée \\
    \bottomrule
  \end{tabularx}
\end{table}

\begin{boiteSynthese}{Synthèse Stratégique Pleine Page (Largeur Utile : 180\,mm)}
L'utilisation de la géométrie large permet de concevoir des blocs récapitulatifs d'une clarté exemplaire :
\begin{itemize}
  \item Les tableaux multi-colonnes disposent de l'amplitude nécessaire pour afficher des phrases complètes
  sans césures disgracieuses ni écrasement typographique.
  \item La transition d'entrée et de sortie s'effectue automatiquement via \code{\textbackslash cleardoublepage},
  garantissant qu'aucune page bâtarde ne présente des marges asymétriques au milieu d'un feuillet.
  \item Le moteur de calcul \LaTeX{} réévalue \code{\textbackslash linewidth} et \code{\textbackslash textwidth}
  à l'intérieur du bloc, assurant la compatibilité totale avec les boîtes \code{tcolorbox}.
\end{itemize}
\end{boiteSynthese}

\end{WideGeometry}
\activerStyleStandard

% ==============================================================================
% 4.3 Démonstration de la Géométrie de Révision (AnnotationGeometry) et ThesisDebug
% ==============================================================================
\section{Démonstration en Direct : Géométrie de Révision et En-Têtes de Débogage}
\label{sec:demo-geometrie-annotation}

Dès la fermeture du bloc \code{WideGeometry},
la commande \code{\textbackslash restoregeometry} rétablit les proportions canoniques de la classe.
Les en-têtes réguliers avec logo sont quant à eux réactivés par la commande
\code{\textbackslash activerStyleStandard}.

Examinons maintenant le scénario inverse : la phase de relecture collaborative,
où l'auteur et les encadrants doivent annoter le manuscrit en vis-à-vis du texte.
en activant le style de diagnostic éditorial \code{\textbackslash activerStyleDebug}.

\begin{AnnotationGeometry}
\activerStyleDebug

\subsection{Marge latérale dédiée de 25~mm et métadonnées d'édition \code{ThesisDebug}}
Observez attentivement la métamorphose de la présente page :
\begin{enumerate}
  \item \textbf{En-tête de contrôle :}
  la mention \code{[\jobname.tex]} apparaît en rouge/primaire à gauche,
  le macaron \textbf{\textcolor{ThesisAccent}{[DEBUG / RÉVISION]}} est mis en exergue au centre,
  et le titre de la section courante s'affiche à droite.
  \item \textbf{Pied de page de traçabilité :}
  le folio affiche désormais la pagination absolue sous la forme \code{Page \thepage{} sur \pageref{LastPage}}.
  \item \textbf{Marge d'annotation :}
  la marge droite (ou extérieure) est configurée à 25\,mm
  via l'assignation \code{marginparwidth=2.5cm},
  offrant l'espace indispensable pour déposer des notes marginales sans empiéter sur le texte.
\end{enumerate}

Voici un exemple de passage soumis à une relecture croisée tripartite :

Dans les architectures documentaires classiques et les chaînes de publication patrimoniales,
\todoH[AuthorStyle]{Préciser la norme ISO relative à la structure des documents XML/PDF}{la rigidité des formats}
constitue souvent le principal obstacle technique au déploiement en entreprise.
Chaque révision éditoriale engendre des régressions de mise en page imprévisibles
qui consomment un temps précieux lors des phases de relecture collégiale.

Face à ces contraintes récurrentes, les équipes hésitent fréquemment entre les traitements de texte WYSIWYG
et les langages de balisage légers incapables de composer des équations différentielles ou des tableaux conformes aux canons des revues scientifiques.
La tentation de recourir à des outils propriétaires génère une dépendance technologique coûteuse
et empêche l'automatisation des tests d'intégrité documentaire au sein d'une forge collaborative.

L'approche modulaire incarnée par notre classe réconcilie ces impératifs contradictoires
\todoH[ReviewerStyle]{Ajouter une référence vers le chapitre 5 pour le pipeline latexmk}{en garantissant une compilation reproductible}
et un découplage total entre le fond, la forme et le profil de diffusion.
Toutes les notes marginales ci-contre comportent une infobulle interactive au survol
et sont automatiquement téléportées vers le \term{Tableau de Bord des Révisions}
accessible en tête de document.

\begin{callout}[Comportement en mode publication (\code{print} ou sans option \code{todos})]
Lorsque l'option de classe \code{todos} est désactivée pour la publication finale,
toutes les boîtes jaunes de marge et les surlignages disparaissent instantanément du texte,
sans laisser la moindre trace ni altérer la justification des lignes.
\end{callout}

\end{AnnotationGeometry}
\activerStyleStandard

% ==============================================================================
% 4.4 Composition Multi-Colonnes Dynamique (multicol)
% ==============================================================================
\section{Composition Multi-Colonnes Dynamique sans Rupture de Page}
\label{sec:demo-multicol}

Dans le standard \LaTeX{}, le passage en deux colonnes s'effectue traditionnellement
par l'instruction primitive \code{\textbackslash twocolumn}.
Cette commande présente toutefois un défaut rédhibitoire en milieu de chapitre :
elle force un saut de page immédiat (\code{\textbackslash clearpage}),
éjecte tous les flottants en attente,
et s'avère incapable d'équilibrer la hauteur des colonnes en bas de section.

Pour surmonter ces écueils, notre configuration intègre le paquet de référence \code{multicol}.
L'environnement \code{multicols} offre trois propriétés exceptionnelles :
\begin{itemize}
  \item Il s'initialise et se clôture n'importe où sur une page en cours sans forcer de saut de page.
  \item Il équilibre automatiquement la hauteur des colonnes sur la dernière page.
  \item Il permet de tracer des filets verticaux séparateurs personnalisables et colorés.
\end{itemize}

\subsection{Démonstration en deux colonnes avec filet vertical séparateur}

Voici une démonstration concrète en deux colonnes avec un filet séparateur
de 0.4\,pt teinté dans la palette secondaire du document :

\setlength{\columnseprule}{0.4pt}
\renewcommand{\columnseprulecolor}{\color{ThesisSecondary!40}}
\setlength{\columnsep}{20pt}

\begin{callout}[Chapeau introductif transversal sur deux colonnes]
Ce bloc d'introduction s'étend sur toute la largeur de la page (\code{\textbackslash textwidth}),
tandis que les paragraphes ci-dessous se distribuent harmonieusement de part et d'autre du filet séparateur.
\end{callout}

\begin{multicols}{2}

\subsubsection{Colonne A : Fondements théoriques}
La distribution du texte sur deux colonnes étroites réduit la longueur de chaque ligne
à environ 35 à 45 caractères.
Cette configuration est universellement adoptée par les revues IEEE, ACM et Elsevier
pour les publications à forte densité de contenu.

Dans ce format, l'algorithme de saut de ligne de Knuth ajuste les espacements inter-mots
avec une sensibilité accrue.
L'activation conjointe du paquet \code{microtype} devient indispensable
pour éviter les lignes lâches (\textit{underfull hbox}) grâce à la saillie optique des marges.

Le passage d'une colonne à l'autre peut être laissé à l'équilibrage automatique,
ou provoqué explicitement à un endroit stratégique grâce à la commande \code{\textbackslash columnbreak}.

\columnbreak

\subsubsection{Colonne B : Implémentation pratique}
Nous forçons ici le passage en colonne B via \code{\textbackslash columnbreak}.
Les équations mathématiques en ligne ou de taille modérée
s'intègrent avec une élégance naturelle :

\begin{equation}
  \label{eq:densite-energie}
  \mathcal{E}_c = \frac{1}{2} \int_{\Omega} \norm{\nabla \vect{u}(\vect{x})}^2 \diff \vect{x}
\end{equation}

Les listes de contrôle et énumérations techniques
demeurent particulièrement compactes et digestes :
\begin{itemize}
  \item filet vertical via \code{\textbackslash columnseprule},
  \item espacement réglé par \code{\textbackslash columnsep},
  \item équilibrage dynamique automatique en bas de page.
\end{itemize}

\end{multicols}

\setlength{\columnseprule}{0pt} % Rétablissement du filet par défaut

\subsection{Démonstration en trois colonnes pour nomenclature technique}

L'environnement \code{multicols} s'étend avec une flexibilité identique à trois colonnes ou plus.
Cette mise en page est idéale pour les dictionnaires de variables,
les tables de nomenclature ou les inventaires de paramètres expérimentaux :

\begin{multicols}{3}
\raggedcolumns
\raggedright
\small

\noindent{\bfseries\sffamily\textcolor{ThesisPrimary}{Paramètres d'état :}}\par\vspace{2pt}
$\vect{x} \in \setR^n$ : Position d'état.\par
$\vect{v} \in \setR^n$ : Vitesse instantanée.\par
$\vect{a} \in \setR^n$ : Accélération.\par
$t \in [0, T]$ : Horizon temporel.\par
$m \in \setR^+$ : Masse effective.

\columnbreak

\noindent{\bfseries\sffamily\textcolor{ThesisPrimary}{Opérateurs tensoriels :}}\par\vspace{2pt}
$\mat{M}$ : Matrice d'inertie.\par
$\mat{C}$ : Forces centrifuges.\par
$\mat{K}$ : Matrice de raideur.\par
$\mat{\Lambda}$ : Tenseur de diffusion.\par
$\mat{I}_n$ : Matrice identité.

\columnbreak

\noindent{\bfseries\sffamily\textcolor{ThesisPrimary}{Critères et Métriques :}}\par\vspace{2pt}
$\mathcal{L}_2$ : Norme quadratique.\par
$\mathrm{SNR}$ : Rapport signal/bruit.\par
$\Delta t$ : Pas d'intégration.\par
$\varepsilon_{\mathrm{tol}}$ : Tolérance de résidu.\par
$\eta$ : Taux de convergence.

\end{multicols}

% ==============================================================================
% 4.5 Rotation Dynamique en Mode Paysage (pdflscape)
% ==============================================================================
\section{Pages Tournées en Mode Paysage avec Orientation PDF (\code{pdflscape})}
\label{sec:demo-landscape}

L'insertion de panoramas expérimentaux, de chronogrammes Gantt
ou de matrices comparatives dépassant 8 à 10 colonnes impose parfois d'orienter le feuillet à 90 degrés.
L'usage historique du paquet standard \code{lscape} pose un problème majeur :
le texte est tourné sur la feuille, mais le lecteur de PDF (Adobe Acrobat, SumatraPDF, aperçu de navigateur)
continue d'afficher la page verticalement, obligeant le lecteur à se contorsionner la tête !

Dans notre atelier, nous utilisons le paquet moderne \code{pdflscape}.
Lors de la compilation avec Lua\LaTeX{}, \code{pdflscape} insère l'attribut PDF interne \code{/Rotate 90} :
la page pivote automatiquement dans le lecteur numérique pour une lecture horizontale immédiate,
tout en préservant le sens d'impression naturel pour le relieur papier.

\begin{landscape}
\activerStyleFocus

\subsection{Matrice synoptique multiparamétrique en orientation paysage}
Observez l'orientation de cette page dans votre visionneuse PDF :
le format A4 est basculé automatiquement en orientation paysage ($297 \times 210$\,mm).
Nous pouvons déployer un tableau synoptique de grande envergure comparant sept moteurs et technologies de composition :

\begin{table}[htbp]
  \centering
  \small
  \caption[Benchmarking technique comparatif des moteurs de composition scientifique]{Benchmarking technique comparatif des technologies de composition scientifique (page paysage dynamique \code{pdflscape})}
  \label{tab:benchmark-paysage}
  \begin{tabularx}{\linewidth}{l c c c c c >{\raggedright\arraybackslash}X}
    \toprule
    \textbf{Technologie} & \textbf{Moteur} & \textbf{Unicode/UTF-8} & \textbf{Microtype} & \textbf{Polices OTF} & \textbf{Vitesse} & \textbf{Verdict en Thèse Doctorale \& Ingénierie} \\
    \midrule
    \textbf{\LaTeXe{} Historique} & pdf\TeX{} & Partiel & Oui & Non (Type 1) & Très rapide & Obsolète pour nouveaux projets ; encodage fragile. \\
    \addlinespace
    \textbf{Xe\LaTeX{}} & Xe\TeX{} & Natif & Partiel & Oui (système) & Rapide & Bon support des fontes, mais microtypographie limitée. \\
    \addlinespace
    \textbf{Lua\LaTeX{}} & Lua\TeX{} & Natif & Complet & Oui (direct/relatif) & Modérée & \textbf{Standard d'excellence recommandé} : programmabilité totale, polices locales figées, UTF-8 robuste. \\
    \addlinespace
    \textbf{Typst} & Rust & Natif & En cours & Oui (système) & Ultra-rapide & Écosystème émergent prometteur, mais support biblio et éditeurs encore jeune pour une thèse d'État. \\
    \addlinespace
    \textbf{MS Word / LibreOffice} & Propriétaire & Natif & Non & Oui & Temps réel & Réprouvé : rupture de reproductibilité, gestion catastrophique des équations et des références croisées. \\
    \bottomrule
  \end{tabularx}
\end{table}

\begin{boiteSynthese}{Règles d'Or pour les Pages en Mode Paysage}
Trois principes garantissent une intégration sans bavure des pages tournées :
\begin{enumerate}
  \item \textbf{Encapsulation stricte :} l'environnement \code{landscape} doit toujours être ouvert
  immédiatement après un saut de page propre et refermé avant de poursuivre la prose courante.
  \item \textbf{Adaptation des en-têtes :} combiner \code{landscape} avec \code{\textbackslash activerStyleFocus}
  permet de masquer les en-têtes horizontaux qui viendraient encombrer inutilement le bord long.
  \item \textbf{Largeur universelle :} utilisez toujours \code{\textbackslash linewidth} dans les environnements \code{tabularx}
  pour occuper exactement la largeur disponible sans calcul manuel hasardeux.
\end{enumerate}
\end{boiteSynthese}

\end{landscape}
\activerStyleStandard

% ==============================================================================
% 4.6 Bonnes Pratiques et Synthèse Opérationnelle
% ==============================================================================
\section{Checklist de l'Architecte : Gestion des Géométries et des En-têtes}
\label{sec:checklist-geometrie}

Pour clore cette démonstration de puissance, le \cref{tab:checklist-geometries} synthétise
les commandes, les contextes d'application et les précautions associées
à chaque modification d'empagement dans notre architecture modulaire.

\begin{table}[htbp]
  \centering
  \small
  \caption[Synthèse opérationnelle des commandes de géométrie et d'en-tête]{Synthèse opérationnelle des commandes de modification d'empagement et d'en-tête}
  \label{tab:checklist-geometries}
  \begin{tabularx}{\linewidth}{>{\raggedright\arraybackslash}p{3.5cm} >{\raggedright\arraybackslash}X >{\raggedright\arraybackslash}X >{\raggedright\arraybackslash}p{2.7cm}}
    \toprule
    \textbf{Commande / Bloc} & \textbf{Effet produit sur le document} & \textbf{Précautions obligatoires} & \textbf{Paquet requis} \\
    \midrule
    \code{\textbackslash geometrieStandard} &
    Restaure la géométrie KOMA-Script (\code{DIV=12}). &
    Exécute automatiquement \code{\textbackslash cleardoublepage}. &
    \code{geometry} \\
    \addlinespace
    \code{WideGeometry} &
    Marges réduites à 15\,mm pour données horizontales denses. &
    Associer à \code{\textbackslash activerStyleFocus} pour éviter les surcharges. &
    \code{config/macros} \\
    \addlinespace
    \code{AnnotationGeometry} &
    Marge latérale élargie à 25\,mm pour les notes de relecture. &
    Idéal en phase de travail avec \code{\textbackslash activerStyleDebug}. &
    \code{config/macros} \\
    \addlinespace
    \code{\textbackslash activerStyle}\allowbreak\code{Focus} &
    En-tête épuré sans filet, folio centré discret. &
    À rétablir via \code{\textbackslash activerStyle}\allowbreak\code{Standard} en fin de bloc. &
    \code{scrlayer-scrpage} \\
    \addlinespace
    \code{\textbackslash activerStyle}\allowbreak\code{Debug} &
    Affiche le fichier source, le statut révision et le folio X/Y. &
    Utilise la référence \code{LastPage} (nécessite deux passes). &
    \code{scrlayer-scrpage} \\
    \addlinespace
    \code{multicols\{N\}} &
    Composition fluide sur $N$ colonnes sans saut de page. &
    Les flottants \code{figure} standards ne sont pas autorisés dans le bloc. &
    \code{multicol} \\
    \addlinespace
    \code{landscape} &
    Rotation physique de la page avec balise PDF \code{/Rotate 90}. &
    Isoler entièrement le tableau entre l'ouverture et la fermeture. &
    \code{pdflscape} \\
    \bottomrule
  \end{tabularx}
\end{table}

Grâce à cette boîte à outils cohérente et modulaire,
l'auteur conserve une liberté éditoriale totale
tout en offrant à son lecteur un manuscrit d'une rigueur typographique sans faille.
